\section{Выводы}

Практический итог работы в том, что исследование качества программы позволило отделить реальные проблемы от предполагаемых. В baseline-версии был найден один действительно важный недочёт: лишние накладные расходы в \texttt{save/load} из-за большого числа мелких файловых операций. После перехода к внутренней буферизации этот участок удалось ускорить без изменения интерфейса словаря и формата файла.

Не менее важно, что исследование не остановилось на первой находке. Отдельно были проверены парсинг команд, рекурсивные функции и работа с памятью. Существенных проблем там не обнаружилось: словарь корректно обрабатывает команды, утечек памяти нет, а глубина рекурсии остаётся безопасной для рассматриваемых размеров дерева.

Для меня эта лабораторная работа оказалась полезной именно как опыт инженерной проверки уже готовой программы. Здесь пришлось не только запускать замеры, но и сопоставлять разные источники информации: профиль времени, нагрузочный бенчмарк и анализ памяти. В результате стало понятнее, как локально улучшать программу и как обосновывать такие изменения не рассуждениями, а измерениями.
\pagebreak
